\color{unidarkgreen}{\subsection{Логика сигнализации}}
\color{black}

\begin{enumerate}[label=\arabic{section}.\arabic{subsection}.\arabic*,  align=left, labelsep=4pt, leftmargin=0pt, itemindent=57pt]

\item
Устройство формирует выходное воздействие на сигнализацию \textit{«Сигн. Сраб.(фикс)»} и \textit{«Откл. защ.
(фикс)»} по схеме «ИЛИ» в следующих случаях:
\begin{itemize}
\item Срабатывание от ДТЗ;
\item Срабатывание МТЗ КО;
\item Срабатывание ТЗНП КО;
\item Срабатывание ТЗОП КО;
\item Срабатывание МТЗ НН;
\item Срабатывание МТЗ ОУ;
\item Срабатывание КИ НН;
\item Отключение от ЗДЗ НН;
\item Отключение от УРОВ ТМПо;
\item Отключение от УРОВ ТМПр;
\item Отключение от внешней РЗ КО, ОУ.
\end{itemize}

\item
Устройство формирует выходное воздействие на сигнализацию \textit{«Сигн. неиспр.»} и \textit{«Сигн. неиспр.
(фикс)»} по схеме «ИЛИ» в следующих случаях:
\begin{itemize}
\item Неисправность цепей тока ДТЗ;
\item Неисправность ОТ (ЗДЗ НН/ТМПо/ТМПр);
\item Потеря GOOSE.
\end{itemize}

\item
Сброс фиксации производится с помощь \textit{«Сброс»}.

\begin{figure}[H]
\centering
\includegraphics[width=0.7\textwidth,height=0.7\textheight,keepaspectratio]{\fbpath/Логика сигнализации/img/LS.pdf}
\captionof{figure}{Функциональный блок логики сигнализации}\label{pict_ls:LS}
\end{figure}

\end{enumerate}